\documentclass[french]{article}

\usepackage[utf8]{inputenc}
\usepackage[T1]{fontenc}
\usepackage{lmodern}
\usepackage{geometry}
\usepackage{graphicx}
\usepackage{babel}
\usepackage{amsmath}
\usepackage{amsthm}
\usepackage{amssymb}
\usepackage{hyperref}
\usepackage{stmaryrd}
\usepackage{enumitem}
\usepackage[most]{tcolorbox}
\usepackage{dsfont}
\usepackage{multirow}
\usepackage{etoc}
\usepackage{titlesec}
\usepackage{esint}
\usepackage{cancel}
\usepackage{mathdots}
\usepackage{indentfirst}

\setlength{\parindent}{0pt}

\setlist[itemize]{label=$\bullet$}
\setlist{before=\leavevmode, leftmargin=*}

\renewcommand{\P}{\mathbb{P}}
\newcommand{\N}{\mathbb{N}}
\newcommand{\Z}{\mathbb{Z}}
\newcommand{\Q}{\mathbb{Q}}
\newcommand{\R}{\mathbb{R}}
\newcommand{\C}{\mathbb{C}}
\newcommand{\K}{\mathbb{K}}
\renewcommand{\L}{\mathbb{L}}
\newcommand{\U}{\mathbb{U}}
\newcommand{\st}{^*}
\newcommand{\tq}{\middle|}
\newcommand{\inv}{^{-1}}
\newcommand{\E}{\mathbb{E}}
\newcommand{\F}{\mathbb{F}}
\newcommand{\V}{\mathbb{V}}
\renewcommand{\ker}{\text{Ker}}
\newcommand{\im}{\text{Im}}
\newcommand{\card}{\text{Card}}
\newcommand{\Tr}{\text{Tr}}
\newcommand{\Vect}{\text{Vect}}
\newcommand{\rg}{\text{rg}}
\newcommand{\vertiii}[1]{{\left\vert\kern-0.25ex\left\vert\kern-0.25ex\left\vert #1 \right\vert\kern-0.25ex\right\vert\kern-0.25ex\right\vert}}
\renewcommand{\o}{\text{o}}
\renewcommand{\O}{\text{O}}
\renewcommand{\d}{\text{d}}
\newcommand{\Id}{\text{Id}}


\theoremstyle{plain}
\newtheorem*{ind}{Indication}

\theoremstyle{definition}

\newtheorem*{ex}{Exemple}
\newtheorem*{rmq}{Remarque}
\newtheorem*{notation}{Notation}
\newtheorem*{rappel}{Rappel}
\newtheorem*{terminologie}{Terminologie}
\newtheorem*{attention}{Attention}
\newtheorem*{conséquence}{Conséquence}

\theoremstyle{remark}
\newtheorem*{app}{Application}

\newtcolorbox[auto counter]{dfn}[1][]{
	enhanced,
	colback=white,
	colframe=black,
	sharp corners,
	boxrule=0.8pt,
	fonttitle=\bfseries,
	colbacktitle=white,
	coltitle=black,
	attach boxed title to top left={yshift=-1mm,xshift=-0.5mm},
	boxed title style={size=minimal, right=2pt, sharp corners, colframe=white},
	title={Définition \thetcbcounter \ifstrempty{#1}{}{ (#1)}},
	before skip=0.5em,
	after skip=0.5em
}

\newtcolorbox[auto counter]{exo}[1][]{
	enhanced,
	arc=1mm,
	colback=white,
	colframe=black,
	coltitle=black,
	fonttitle=\bfseries,
	boxrule=0.8pt,
	shadow={1.2pt}{-1.2pt}{0pt}{black!50},
	attach title to upper={\ },
	title={Exercice \thetcbcounter\ifblank{#1}{}{ #1}},
	before skip=0.5em,
	after skip=0.5em,
	left=2mm, right=2mm, top=1mm, bottom=1mm
}

\newtcolorbox{met}[1][]{
	enhanced,
	arc=1mm,
	colback=white,
	colframe=black,
	coltitle=black,
	fonttitle=\bfseries,
	boxrule=0.8pt,
	shadow={1.2pt}{-1.2pt}{0pt}{black!50},
	attach title to upper={\ },
	title={Point méthode \ifblank{#1}{}{#1}},
	before skip=1em,
	after skip=1em,
	left=2mm, right=2mm, top=1mm, bottom=1mm
}

\newcounter{thmcount}

\newtcolorbox[use counter=thmcount]{thm}[1][]{
	enhanced,
	colback=white,
	colframe=black,
	sharp corners,
	left skip=0mm,
	boxrule=1.3pt,
	fonttitle=\bfseries,
	colbacktitle=white,
	coltitle=black,
	attach boxed title to top left={yshift=-1mm,xshift=-0.5mm},
	boxed title style={size=minimal, right=2pt, sharp corners, colframe=white},
	title={Théorème \thethmcount \ifstrempty{#1}{}{ (#1)}},
	before skip=0.5em,
	after skip=0.5em
}

\newtcolorbox[use counter=thmcount]{prop}[1][]{
	enhanced,
	colback=white,
	colframe=black,
	sharp corners,
	boxrule=0.8pt,
	fonttitle=\bfseries,
	colbacktitle=white,
	coltitle=black,
	attach boxed title to top left={yshift=-1mm,xshift=-0.5mm},
	boxed title style={size=minimal, right=2pt, sharp corners, colframe=white},
	title={Proposition \thethmcount \ifstrempty{#1}{}{ (#1)}},
	before skip=0.5em,
	after skip=0.5em
}

\newtcolorbox[use counter=thmcount]{cor}[1][]{
	enhanced,
	colback=white,
	colframe=black,
	sharp corners,
	boxrule=0.8pt,
	fonttitle=\bfseries,
	colbacktitle=white,
	coltitle=black,
	attach boxed title to top left={yshift=-1mm,xshift=-0.5mm},
	boxed title style={size=minimal, right=2pt, sharp corners, colframe=white},
	title={Corollaire \thethmcount \ifstrempty{#1}{}{ (#1)}},
	before skip=0.5em,
	after skip=0.5em
}

\newtcolorbox[use counter=thmcount]{lem}[1][]{
	enhanced,
	colback=white,
	colframe=black,
	sharp corners,
	boxrule=0.5pt,
	fonttitle=\bfseries,
	colbacktitle=white,
	coltitle=black,
	attach boxed title to top left={yshift=-1mm,xshift=-0.5mm},
	boxed title style={size=minimal, right=2pt, sharp corners, colframe=white},
	title={Lemme \thethmcount \ifstrempty{#1}{}{ (#1)}},
	before skip=0.5em,
	after skip=0.5em
}

\title{Endomorphismes d'un espace euclidien}
\date{}


\begin{document}
\tableofcontents

\newpage
\maketitle
Dans ce chapitre, on considère un espace euclidien $E$ dont le produit scalaire sera noté $(\cdot\mid\cdot)$. La norme et la distance euclidiennes associées seront respectivement notées $\lVert\cdot\rVert$ et $d(\cdot,\cdot)$.

\section{Rappels, compléments}

\subsection{Formes bilinéaires symétriques, formes quadratiques (HP)}

\begin{dfn}
	On dit qu'une forme bilinéaire $\varphi$ sur $E$ est symétrique lorsque $$\forall (x,y)\in E^2,\ \varphi(x,y)=\varphi(y,x).$$ On note $\mathcal{BLS}(E)$ l'ensemble des formes bilinéaires symétriques sur $E$.
\end{dfn}

\begin{dfn}
	A toute forme bilinéaire symétrique $\varphi$ sur $E$ on associe l'application $q$ de $E$ dans $\R$ définie par $$\forall x\in E,\ q(x)=\varphi(x,x)$$ C'est la forme quadratique associée à $\varphi$. Réciproquement, une forme quadratique sur $E$ est une application $q:E\to\R$ pour laquelle il existe une forme bilinéaire symétrique $\varphi$ vérifiant $$\forall x\in E,\ q(x)=\varphi(x,x)$$ On note $\mathcal{Q}(E)$ l'ensemble des formes quadratiques.
\end{dfn}

\begin{rmq}
	Les formules de polarisation permettent de retrouver la forme bilinéaire symétrique $\varphi$ à partir de la forme quadratique $q$ : $$\forall (x,y)\in E^2,\ \varphi(x,y)=\frac{q(x+y)-q(x)-q(y)}{2}=\frac{q(x)+q(y)-q(x-y)}{2}=\frac{q(x+y)-q(x-y)}{4}.$$
\end{rmq}

\subsubsection*{Matrice d'une forme bilinéaire}
Soit $\mathcal{B}=(e_1,\ldots,e_n)$ une base de $E$. A toute forme bilinéaire $\varphi$ sur $E$, on associe sa matrice dans $\mathcal{B}$ définie par : $$M_\mathcal{B}(\varphi)=\left(\varphi(e_i,e_j)\right)_{1\leqslant i,j\leqslant n}$$ Cette matrice est symétrique si, et seulement si, $\varphi$ est une forme bilinéaire symétrique. De plus, si $(x,y)\in E^2$ et si on pose $X=M_\mathcal{B}(x)$ et $Y=M_\mathcal{B}(y)$, alors on dispose de la relation : $$\varphi(x,y)=X^TM_\mathcal{B}(\varphi)Y.$$

\begin{exo}
	Effet d'un changement de base sur la matrice d'une forme bilinéaire.
\end{exo}

\begin{dfn}
	Une forme bilinéaire symétrique $\varphi$ (ou sa forme quadratique associée) est dite positive lorsque $$\forall x\in E,\ \varphi(x,x)\geqslant 0$$ Elle est définie positive si de plus $$\forall x\in E,\ \varphi(x,x)=0\implies x=0.$$
\end{dfn}

\begin{dfn}
	Un produit scalaire est une forme bilinéaire symétrique définie positive.
\end{dfn}

\begin{ex}
	\begin{enumerate}
		\item L'application : $$\begin{array}{rcl}
			\R^n\times\R^n&\longrightarrow&\R \\
			(x,y)&\longmapsto&\sum_{i=1}^{n}x_iy_i
		\end{array}$$ est un produit scalaire sur $\R^n$, appelé \textbf{produit scalaire canonique de $\R^n$}.
		\item Plus généralement, si $E$ est muni d'une base $(e_i)_{i\in I}$, on peut définir un produit scalaire sur $E$ en posant : $$\forall (x_i)\in\K^{(I)},\ \forall (y_i)\in\K^{(I)},\ \left(\sum_{i\in I}x_ie_i\bigg| \sum_{i\in I} y_ie_i\right)=\sum_{i\in I}x_iy_i$$
		\item L'ensemble des suites réelles de carré sommable, c'est-à-dire des suites $(u_n)_{n\in\N}$ réelles telles que $\sum u_n^2$ converge est un sous-espace vectoriel de $\R^\N$ sur lequel $(u,v)\mapsto\displaystyle\sum_{n\in\N}u_nv_n$ définit un produit scalaire.
		\item Soit $[a,b]$ un segment de $\R$ d'intérieur non vide. L'application : $$\varphi:(f,g)\mapsto\int_{a}^{b}f(x)g(x)\ dx$$ définit une forme bilinéaire symétrique positive sur $\mathcal{CM}([a,b],\R)$. C'est un produit scalaire sur $E=\mathcal{C}^0([a,b],\R)$.
	\end{enumerate}
\end{ex}

\begin{dfn}
	On appelle espace préhilbertien réel tout $\R$-espace vectoriel muni d'un produit scalaire. Lorsque l'espace vectoriel est de dimension finie, on parle d'espace vectoriel euclidien.
\end{dfn}

\subsection{Norme associée à un produit scalaire}

On considère $E$ un espace préhilbertien réel.

\begin{dfn}
	On appelle \textbf{norme euclidienne} associée au produit scalaire $(\cdot\mid\cdot)$ l'application : $$\begin{array}{rcl}
		E&\longrightarrow&\R_+ \\
		x&\longmapsto&\lVert x\rVert=\sqrt{(x\mid x)}
	\end{array}$$ et \textbf{distance euclidienne} associée au produit scalaire $(\cdot\mid\cdot)$ l'application : $$\begin{array}{rcl}
	E^2&\longrightarrow&\R_+ \\
	(x,y)&\longmapsto&\lVert x-y\rVert
	\end{array}$$
\end{dfn}

\begin{prop}[Inégalité de Cauchy-Schwarz]
	La norme euclidienne associée au produit scalaire $(\cdot\mid\cdot)$ vérifie : 
	\begin{enumerate}
		\item $\forall (x,y)\in E^2,\ \lvert (x\mid y)\rvert\leqslant\lVert x\rVert\lVert y\rVert$
		\item Cette inégalité est une égalité si, et seulement si, $x$ et $y$ sont colinéaires, c'est-à-dire si, et seulement si : $$\exists\lambda\in\R\ : \ y=\lambda x\ \ \text{ou}\ \ \exists\lambda\in\R\ : \ x=\lambda y$$
	\end{enumerate}
\end{prop}

\begin{ex}
	\begin{enumerate}
		\item Soit $(x_i)_{1\leqslant i\leqslant n}$ et $(y_i)_{1\leqslant i\leqslant n}$ deux vecteurs de $\R^n$. L'inégalité de Cauchy-Schwarz pour le produit scalaire canonique de $\R^n$ s'écrit : $$\left\lvert\sum_{i=1}^{n}x_iy_i\right\rvert\leqslant\left(\sum_{i=1}^{n}x_i^2\right)^{1/2}\left(\sum_{i=1}^{n}y_i^2\right)^{1/2}$$
		\item L'inégalité de Cauchy-Schwarz pour le produit scalaire sur $\mathcal{C}^0([a,b],\R)$ défini par $$(f,g)\mapsto\int_{a}^{b}fg$$ s'écrit : $$\left\lvert \int_{a}^{b}fg\right\rvert\leqslant\left(\int_{a}^{b}f^2\right)^{1/2}\left(\int_{a}^{b}g^2\right)^{1/2}$$
	\end{enumerate}
\end{ex}

\begin{prop}
	La norme euclidienne $\lVert\cdot\rVert$ associée au produit scalaire $(\cdot\mid\cdot)$ est une norme. Elle vérifie donc en particulier l'inégalité triangulaire : $$\forall (x,y)\in E^2,\ \lVert x+y\rVert\leqslant\lVert x\rVert+\Vert y\rVert$$ En plus de l'inégalité triangulaire, on dispose dans le cas d'une norme euclidienne d'un cas d'égalité : il y a égalité dans l'inégalité triangulaire si, et seulement si, $x$ et $y$ sont colinéaires directs, c'est-à-dire si, et seulement si : $$\exists\lambda\in\R_+\ : \ y=\lambda x\ \ \text{ou}\ \ \exists\lambda\in\R_+\ : \ x=\lambda y$$
\end{prop}

\begin{rmq}[HP]
	Si $\varphi$ est une forme bilinéaire symétrique positive sur $E$, on dispose toujours de l'inégalité de Cauchy-Schwarz : $$\forall (x,y)\in E^2,\ \varphi(x,y)^2\leqslant\varphi(x,x)\varphi(y,y)$$ On peut alors poser $N(x)=\sqrt{\varphi(x,x)}$ pour tout $x\in E$, définissant ainsi seulement une semi-norme sur $E$. La condition d'égalité dans l'inégalité de Cauchy-Schwarz devient : $$\exists\lambda\in\R\ : \ N(x+\lambda y)=0\ \ \text{ou}\ \ \exists\lambda\in\R\ : \ N(y+\lambda x)=0$$ et de même pour la condition d'égalité dans l'inégalité triangulaire.
\end{rmq}

\begin{ex}
	Si $f$ et $g$ sont des fonctions réelles continues par morceaux sur $[a,b]$ avec $a<b$, alors : $$\left(\int_{a}^{b}fg\right)^2\leqslant\left(\int_{a}^{b}f\right)\left(\int_{a}^{b}g\right)$$ avec égalité si, et seulement si, les restrictions de $f$ et $g$ au complémentaire d'une partie finie de $[a,b]$ sont proportionnelles.
\end{ex}

\subsection{Formules de polarisation}

La norme euclidienne d'un espace vectoriel préhilbertien est définie à partir du produit scalaire. Réciproquement, si l'on connaît la norme euclidienne, on peut retrouver le produit scalaire grâce aux égalités suivantes appelées \textbf{formules (ou identités) de polarisation}.

\begin{prop}[Formules de polarisation]
	Pour tout $(x,y)\in E^2$, on a :
	\begin{itemize}
		\item $2(x\mid y)=\lVert x+y\rVert^2-\lVert x\rVert^2-\lVert y\rVert^2=\lVert x\rVert^2+\lVert y\rVert^2-\lVert x-y\rVert^2$ ;
		\item $4(x\mid y)=\lVert x+y\rVert^2-\lVert x-y\rVert^2$.
	\end{itemize}
\end{prop}

En combinant les deux premières formules de polarisation, on obtient le résultat suivant : 

\begin{prop}[Identité du parallélogramme]
	Pour tout $(x,y)\in E^2$, on a $\lVert x+y\rVert^2+\lVert x-y\rVert^2=2(\lVert x\rVert^2+\lVert y\rVert^2)$.
\end{prop}

\begin{rmq}
	Cette égalité traduit le fait que, dans un parallélogramme, la somme des carrés des longueurs des deux diagonales est égale à la somme des carrés des longueurs des quatre côtés.
\end{rmq}

\begin{exo}[$\star$]
	Montrer que réciproquement, toute norme vérifiant l'identité du parallélogramme est issue d'un produit scalaire.
\end{exo}

\subsection{Orthogonalité}
On considère $E$ un espace vectoriel réel muni d'une forme bilinéaire symétrique $\varphi$ (en pratique, d'un produit scalaire).
\subsubsection*{Définitions}

\begin{dfn}
	\begin{itemize}
		\item On dit que deux vecteurs $x$ et $y$ de $E$ sont \textbf{orthogonaux} et l'on note $x\bot y$ lorsque $\varphi(x,y)=0$.
		\item Si $A$ et $B$ sont deux parties de $E$, on dit que $A$ est orthogonale à $B$, et l'on note $A\bot B$, lorsque : $$\forall (x,y)\in A\times B,\ \varphi(x,y)=0$$
		\item Si $A$ est une partie de $E$, l'\textbf{orthogonal} de $A$ est l'ensemble des vecteurs orthogonaux à tous les éléments de $A$. on le note $A^\bot$.
	\end{itemize}
\end{dfn}

\begin{exo}
	Montrer : 
	\begin{enumerate}
		\item $A\bot A^\bot$ ;
		\item $B\subset A^\bot \iff A\bot B\iff A\subset B^\bot$.
	\end{enumerate}
\end{exo}

\subsubsection*{Résultats}

\begin{enumerate}
	\item $A^\bot$ est un sous-espace vectoriel de $E$.
	\item $A\subset B\implies B^\bot\subset A^\bot$.
	\item $\text{Vect}(A)^\bot=A^\bot$.
	\item $A\subset (A^\bot)^\bot$.
	\item Pour un produit scalaire : $x\in A\cap A^\bot\implies x=0$.
\end{enumerate}

\subsubsection*{Bases orthogonales}

\begin{dfn}
	Une \textbf{famille orthogonale} est une famille de vecteurs deux à deux orthogonaux pour $\varphi$.
\end{dfn}

\begin{rmq}
	\begin{itemize}
		\item On dit qu'un vecteur $x$ est isotrope pour $\varphi$ lorsque $\varphi(x,x)=0$.
		\item Seul $0$ est isotrope dans un espace préhilbertien.
		\item Toute famille orthogonale constituée de vecteurs non isotropes est libre.
		\item En particulier, dans un espace préhilbertien, toute famille orthogonale constituée de vecteurs non nuls est libre.
	\end{itemize}
\end{rmq}

\begin{exo}
	\begin{enumerate}
		\item L'ensemble des vecteurs isotropes est-il un sous-espace vectoriel ?
		\item Montrer que c'est le cas pour une forme quadratique positive.
	\end{enumerate}
\end{exo}

\begin{thm}
	Si $E$ est de dimension finie, il admet une base orthogonale.
\end{thm}

\subsection{Bases orthonormées}

Soit $E$ un espace préhilbertien réel.

\begin{dfn}
	\begin{itemize}
		\item On appelle vecteur normé, ou unitaire, tout vecteur de norme $1$.
		\item On appelle \textbf{famille orthogonale} de $E$ toute famille de vecteurs de $E$ deux à deux orthogonaux.
		\item On appelle \textbf{famille orthonormée} (ou \textbf{orthonormale}) de $E$ tout famille de vecteurs de $E$ normés et deux à deux orthogonaux.
	\end{itemize}
\end{dfn}

\begin{prop}
	Tout espace vectoriel euclidien admet une base orthonormée.
\end{prop}

\begin{ex}
	\begin{enumerate}
		\item Dans $\R^n$ muni du produit scalaire canonique, les vecteurs de la base canonique sont normés et orthogonaux deux à deux.
		\item Dans l'espace vectoriel des fonctions continues et $2\pi$-périodiques sur $\R$, muni du produit scalaire : $$(f,g)\mapsto\frac{1}{\pi}\int_{0}^{2\pi}f(x)g(x)\ dx$$ les fonctions $(f_k:x\mapsto\cos(kx))_{k\in\N}$ et $(g_k:x\mapsto\sin(kx))_{k\in\N}$ forment une famille orthogonale.
	\end{enumerate}
\end{ex}

\begin{exo}
	Dans le deuxième exemple ci-dessus, s'agit-il d'une famille orthonormée ?
\end{exo}

\begin{prop}
	Soit $E$ un espace vectoriel euclidien et $\mathcal{B}=(e_1,\ldots,e_n)$ une base orthonormée de $E$.
	\begin{enumerate}
		\item Si $x$ est un vecteur de $E$, alors on a $x=\displaystyle\sum_{i=1}^{n}(e_i\mid x)e_i$.
		\item Si $x=\displaystyle\sum_{i=1}^{n}x_ie_i$ et $y=\displaystyle\sum_{i=1}^{n}y_ie_i$ sont deux vecteurs de $E$, alors on a : $$(x\mid y)=\sum_{i=1}^{n}x_iy_i=X^TY\ \ \text{et}\ \ \lVert x\rVert^2=\sum_{i=1}^{n}x_i^2=X^TX$$ où $X$ et $Y$ sont les matrices colonnes constituées des composantes dans la base $\mathcal{B}$ des vecteurs $x$ et $y$.
	\end{enumerate}
\end{prop}

\begin{exo}
	Soit $O(t)$ un point et $\overrightarrow{u}(t)$, $\overrightarrow{v}(t)$ deux vecteurs orthonormées de l'espace, et $M(t)$ un point du plan $(O(t),\overrightarrow{u}(t),\overrightarrow{v}(t))$. On suppose que les fonctions $O$, $\overrightarrow{u}$, $\overrightarrow{v}$ et $M$ sont continues. Montrer que les coordonnées de $M(t)$ dans $(O(t),\overrightarrow{u}(t),\overrightarrow{v}(t))$ sont des fonctions continues de $t$.
\end{exo}

\subsection{Projection orthogonale}

Soit $E$ un espace vectoriel préhilbertien.

\subsubsection*{Orthogonal d'un sous-espace de dimension finie}

\begin{thm}
	Pour tout sous-espace vectoriel $F$ de dimension finie de $E$, on a $E=F\oplus F^\bot$. La \textbf{projection orthogonale} sur $F$, notée $\pi_F$ est la projection sur $F$ parallèlement à son supplémentaire orthogonal $F^\bot$. Si $\mathcal{B}=(e_1,\ldots,e_p)$ est une base orthonormée de $F$, le projeté orthogonal sur $F$ d'un vecteur $x$ de $E$ est $$\pi_F(x)=\sum_{i=1}^{p}(e_i\mid x)e_i$$
\end{thm}

\begin{prop}
	Si $F$ est un sous-espace vectoriel de $E$ et si $E$ est de dimension finie, alors :
	\begin{enumerate}
		\item $\dim(F)+\dim(F^\bot)=\dim(E)$ ;
		\item $(F^\bot)^\bot=F$.
	\end{enumerate}
\end{prop}

\begin{ex}
	En dimension finie, le supplémentaire orthogonal d'un hyperplan est donc une droite, et le supplémentaire orthogonal d'une droite, un hyperplan.
\end{ex}

\subsubsection*{Distance à un sous-espace vectoriel}

\begin{dfn}
	Soit $\mathcal{X}$ une partie non vide de $E$ et $A$ un point de $E$. On appelle \textbf{distance} de $A$ à $\mathcal{X}$ la quantité : $$d(A,\mathcal{X})=\underset{M\in\mathcal{X}}{\inf}d(A,M)$$
\end{dfn}

\begin{thm}[Projection orthogonale sur un sous-espace de dimension finie]
	Soit $F$ un sous-espace vectoriel de dimension finie d'un espace préhilbertien réel $E$ et $x$ un vecteur de $E$. La distance du vecteur $x$ à $F$ est atteinte en un unique point de $F$, à savoir $\pi_F(x)$. Autrement dit : 
	\begin{enumerate}
		\item $d(x,F)=\lVert x-\pi_F(x)\rVert$ ;
		\item $\forall y\in F,\ d(x,F)=\lVert x-y\rVert\iff y=\pi_F(x)$.
	\end{enumerate}
\end{thm}

\begin{ex}
	Calcul de $\displaystyle\underset{(a,b)\in\R^2}{\inf}\int_{0}^{2\pi}(x-a\cos(x)-b\sin(x))^2\ dx$
\end{ex}

\begin{exo}
	Soit $F$ un sous-espace vectoriel de $E$.
	\begin{enumerate}
		\item Soit $x\in E$. Montrer que la distance de $x$ à $F$ est atteinte si, et seulement si, $x\in F+F^\bot$.
		\item Retrouver que si $F$ est de dimension finie, alors $E=F\oplus F^\bot$.
	\end{enumerate}
\end{exo}

\subsection{Procédé d'orthonormalisation de Gram-Schmidt}

\begin{thm}[Algorithme d'orthogonalisation de Gram-Schmidt]
	Soit $\mathcal{F}=(e_1,\ldots,e_n)$ une famille libre de $E$. Alors il existe une famille orthogonale $(g_1,\ldots,g_n)$ de vecteurs (non nuls) de $E$ telle que $$\forall p\in\llbracket 1,n\rrbracket,\ \text{Vect}\{e_1,\ldots,e_p\}=\text{Vect}\{g_1,\ldots,g_p\}$$
\end{thm}

\begin{rmq}
	\begin{itemize}
		\item Si $(e_1,\ldots,e_n)$ est libre, il existe une famille orthonormée $(f_1,\ldots,f_n)$ de $E$ telle que $$\forall p\in\llbracket 1,n\rrbracket,\ \text{Vect}\{e_1,\ldots,e_p\}=\text{Vect}\{f_1,\ldots,f_p\}$$
		\item Dans le procédé d'orthonormalisation de Gram-Schmidt, on construit le vecteur : $$g_{p+1}=e_{p+1}-\sum_{i=1}^{p}\lambda_ig_i$$ et l'on choisit $\lambda_i$ pour que $g_{p+1}-e_{p+1}$ soit orthogonal à $g_i$. On obtient donc le vecteur $g_{p+1}$ en retranchant à $e_{p+1}$ son projeté orthogonal sur $F=\text{Vect}(e_1,\ldots,e_p)$. Pour normaliser le vecteur $g_{p+1}$, il suffit alors d'appliquer le théorème de Pythagore aux vecteurs orthogonaux $\pi_F(e_{p+1})$ et $g_{p+1}$. On obtient ainsi : $$\lVert g_{p+1}\rVert^2=\lVert e_{p+1}\rVert^2-\lVert\pi_F(e_{p+1})\rVert^2=\lVert e_{p+1}\rVert^2-\sum_{i=1}^{p}\lambda_i\lVert g_i\rVert^2$$
	\end{itemize}
\end{rmq}

\begin{exo}
	Montrer que si $(e_1,\ldots,e_n)$ est une famille libre de $E$, alors il existe une unique famille orthonormée $(f_1,\ldots,f_n)$ de $E$ telle que pour tout $p\in\llbracket1,n\rrbracket$, on ait : $$\text{Vect}\{e_1,\ldots,e_p\}=\text{Vect}\{f_1,\ldots,f_p\}\ \ \text{et}\ \ (e_p\mid f_p)>0$$
\end{exo}

\begin{prop}
	Soit $(e_n)_{n\in\N}$ une famille libre d'un espace préhilbertien réel $E$. Alors, il existe une famille orthonormée $(f_n)_{n\in\N}$ telle que $$\forall n\in\N,\ \text{Vect}\{e_0,\ldots,e_n\}=\text{Vect}\{f_0,\ldots,f_n\}$$
\end{prop}

\begin{prop}[HP]
	Tout espace préhilbertien réel de dimension dénombrable possède une base orthonormée.
\end{prop}

\subsection{Matrices de Gram (HP)}

Soit $E$ un espace préhilbertien réel. Pour tout famille finie $(x_i)_{1\leqslant i\leqslant p}$, on pose : $$G(x_1,\ldots,x_p)=((x_i|x_j))_{1\leqslant i,j\leqslant p}\ \ \text{et} \ \ \text{Gram}(x_1,\ldots,x_p)=\det G(x_1,\ldots,x_p)$$
\begin{enumerate}
	\item Montrer que $(x_i)$ est liée si, et seulement si, $\text{Gram}(x_i)=0$.
	\item Montrer qu'ajouter à l'un des $x_i$ une combinaison linéaire des autres ne change pas le rang de $G(x_i)$.
	\item Montrer que $\text{rg}G(x_i)=\text{rg}(x_i)$.
	\item On suppose $E$ muni d'une base $\mathcal{B}=(e_1,\ldots,e_p)$. Donner une relation entre $G(x_i)$, $G(e_i)$ et $M_\mathcal{B}(x_i)$.
	\item Montrer que $G(x_i)$ est positive et en déduire que $\text{Gram}(x_i)\geqslant0$.
	\item Si $F$ est un sous-espace vectoriel dont $y_1$, ..., $y_p$ est une base, on a la formule : $$d(x,F)^2=\frac{\text{Gram}(x,y_1,\ldots,y_p)}{\text{Gram}(y_1,\ldots,y_p)}$$
	\item Lorsque $E$ est de dimension finie $n$, exprimer $\text{Gram}(x_i)$ en fonction du produit mixte des vecteurs $x_i$ pour une orientation quelconque de $E$.
\end{enumerate}

\section{Adjoint d'un endomorphisme}

Dans toute cette section, $E$ est un espace euclidien.

\subsection{Définition}

Pour tout vecteur $a$, l'application $x\mapsto (a\mid x)$ est une forme linéaire sur $E$. Le théorème de Riesz permet de montrer réciproquement que toute forme linéaire sur $E$ est de cette forme.

\begin{thm}[Représentation des formes linéaires sur $E$]
	Soit $\varphi$ une forme linéaire sur $E$. Il existe un unique vecteur $a\in E$ tel que $$\forall x\in E,\ \varphi(x)=(a\mid x)$$
\end{thm}

Le théorème de Riesz permet de définir l'adjoint d'un endomorphisme de $E$.

\begin{thm}[Théorème de représentation de Riesz]
	Soit $u\in\mathcal{L}(E)$. Il existe un unique endomorphisme de $E$, noté $u\st$, tel que $$\forall (x,y)\in E^2,\ (x\mid u(y))=(u\st(x)\mid y)$$ On l'appelle l'\textbf{adjoint} de $u$.
\end{thm}

\begin{rmq}
	La preuve que nous venons de donner de la linéarité de l'adjoint découle d'un résultat plus général : si $u$ et $v$ sont des applications de $E$ dans $E$ vérifiant $$\forall (x,y)\in E^2,\ (x\mid u(y))=(v(x)\mid y)$$ alors $u$ et $v$ sont linéaires.
\end{rmq}

L'adjoint peut aussi exister en dimension infinie.
\begin{ex}
	On munit $\R[X]$ du produit scalaire défini par $$(P\mid Q)=\int_{-\infty}^{+\infty}P(x)Q(x)e^{-x^2}\ dx$$ En notant $D$ l'opérateur de dérivation et $M$ celui de multiplication par $X$, on a $M\st=M$ et $D\st=2M-D$.
\end{ex}

\subsection{Propriétés}

\begin{prop}
	\begin{itemize}
		\item L'application $u\mapsto u\st$ est linéaire sur $\mathcal{L}(E)$.
		\item Pour tout $(u,v)\in \mathcal{L}(E)^2$, dispose des relations : $$(v\circ u)\st=u\st\circ v\st \ \ \text{et}\ \ (u\st)\st=u$$
	\end{itemize}
\end{prop}

\begin{rmq}
	De la propriété précédente découle le fait que $u\mapsto u\st$ est une symétrie de l'espace vectoriel $\mathcal{L}(E)$.
\end{rmq}

\begin{ex}
	\begin{enumerate}
		\item Si $u$ est un automorphisme, alors $u\st$ est un automorphisme et $(u\st)\inv=(u\inv)\st$.
		\item En reprenant l'exemple précédent, l'endomorphisme $D^2-2MD$ est égal à son adjoint.
	\end{enumerate}
\end{ex}

\begin{prop}
	Soit $u\in\mathcal{L}(E)$ et $\mathcal{B}$ une base orthonormée de $E$. Alors $M_\mathcal{B}(u\st)=M_\mathcal{B}(u)^T$.
\end{prop}

\begin{rmq}
	On munit $\R^n$ de sa structure euclidienne canonique. Soit $A\in\mathcal{M}_n(\R)$ dont on note $\varphi_A$ l'endomorphisme canoniquement associé. Alors $(\varphi_A)\st$ est l'endomorphisme canoniquement associé à $A^T$ car la base canonique de $\R^n$ est une base orthonormée.
\end{rmq}

\begin{cor}
	Soit $u\in\mathcal{L}(E)$. Alors on a : $$\det(u\st)=\det(u)\ ; \ \Tr(u\st)=\Tr(u)\ \text{et}\ \text{rg}(u\st)=\text{rg}(u)$$
\end{cor}

\begin{prop}
	Soit $u\in\mathcal{L}(E)$. Alors on a : $$\ker(u\st)=\text{Im}(u)^\bot \ \ \text{et} \ \ \text{Im}(u\st)=\ker(u)^\bot$$
\end{prop}

\begin{rmq}
	En remplaçant $u$ par $u\st$ et en utilisant le fait que $(u\st)\st=u$, on obtient aussi : $$\ker(u)=\text{Im}(u\st)^\bot \ \ \text{et} \ \ \text{Im}(u)=\ker(u\st)^\bot$$
\end{rmq}

\begin{lem}
	Soit $u\in\mathcal{L}(E)$ et $F$ un sous-espace vectoriel de $E$ stable par $u$. Alors $F^\bot$ est stable par $u$.
\end{lem}

\section{Matrices orthogonales}

\subsection{Définition}

\begin{dfn}
	Soit $A\in\mathcal{M}_n(\R)$. On dit que $A$ est une matrice orthogonale lorsque $A^TA=I_n$.
\end{dfn}

\begin{notation}
	On note $\mathcal{O}_n(\R)$ ou encore $\mathcal{O}(n)$ l'ensemble des matrices orthogonales de $\mathcal{M}_n(\R)$.
\end{notation}

\begin{prop}
	Soit $A\in\mathcal{M}_n(\R)$ dont note note $C_1,\ldots,C_n$ les colonnes et $L_1,\ldots,L_n$ les lignes (vues comme des vecteurs de $\R^n$). Alors la matrice $A$ est orthogonale si, et seulement si, l'une des propriétés équivalentes suivantes est vérifiée :
	\begin{enumerate}[label=(\roman*)]
		\item $A^TA=I_n$ ;
		\item $AA^T=I_n$ ;
		\item $A\in\mathcal{GL}_n(\R)$ et $A\inv=A^T$ ;
		\item la famille $(C_1,\ldots,C_n)$ est une base orthonormée ;
		\item la famille $(L_1,\ldots,L_n)$ est une base orthonormée.
	\end{enumerate}
\end{prop}

\begin{ex}
	\begin{enumerate}
		\item La matrice $A=\displaystyle\begin{pmatrix}
			1&1\\
			0&1
		\end{pmatrix}$ n'est pas une matrice orthogonale. En effet, sa deuxième colonne n'est pas unitaire.
		\item Pour tout $\theta\in\R$, la matrice $R(\theta)=\displaystyle\begin{pmatrix}
			\cos(\theta)&-\sin(\theta)\\
			\sin(\theta)&\cos(\theta)
		\end{pmatrix}$ est une matrice orthogonale. En effet, les deux colonnes sont évidemment orthogonales et aussi unitaires car $\cos^2(\theta)+\sin^2(\theta)=1$.
		\item La base canonique de $\R^n$ étant une base orthonormée, tout matrice obtenue en permutant les colonnes et/ou les lignes de la matrice $I_n$ est orthogonale.
		\item Si l'un des coefficients d'une matrice orthogonale vaut $-1$ ou $1$, alors tous les autres coefficients de la ligne et de la colonne sont nuls.
	\end{enumerate}
\end{ex}

La propriété suivante exprime le fait que les matrices de $\mathcal{O}_n(\R)$ sont les matrices de changement de base orthonormée dans n'importe quel espace euclidien de dimension $n$.

\begin{prop}
	Soit $\mathcal{B}$ une base orthonormée de $E$ et $\mathcal{B}'$ une famille de $n$ vecteurs de $E$. Alors $\mathcal{B}'$ est une base orthonormée de $E$ si, et seulement si, $M_\mathcal{B}(\mathcal{B}')\in\mathcal{O}_n(\R)$.
\end{prop}

\begin{dfn}
	Deux matrices $A$ et $B$ de $\mathcal{M}_n(\R)$ sont dites \textbf{orthogonalement semblables} lorsque : $$\exists P\in\mathcal{O}_n(\R)\ : \ A=PBP\inv$$
\end{dfn}

\begin{rmq}
	Deux matrices qui représentent un même endomorphisme dans deux bases orthonormées sont donc orthogonalement semblables en vertu des formules de changement de base et de la propriété précédente.
\end{rmq}

\begin{met}
	Pour montrer que deux matrices $A$ et $B$ de $\mathcal{M}_n(\R)$ sont orthogonalement semblables, on pourra montrer qu'il existe une base orthonormée de $\R^n$ dans laquelle l'endomorphisme canoniquement associé à $A$ a pour matrice $B$.
\end{met}

\subsection{Le groupe orthogonal}

\begin{rmq}
	Soit $A\in\mathcal{O}_n(\R)$. Alors $A^TA=I_n$ puis, par passage au déterminant et par conservation du déterminant par transposition, on obtient $\det(A)^2=1$ donc $\det(A)\in\{-1,1\}$, ce qui justifie la définition suivante.
\end{rmq}

\begin{dfn}
	Une matrice $A\in\mathcal{O}_n(\R)$ est dite :
	\begin{itemize}
		\item \textbf{positive}, ou \textbf{directe}, lorsque $\det(A)=1$ ;
		\item \textbf{négative}, ou \textbf{indirecte}, lorsque $\det(A)=-1$.
	\end{itemize}
\end{dfn}

\begin{notation}
	On note $\mathcal{SO}_n(\R)$ l'ensemble des matrices orthogonales positives de $\mathcal{M}_n(\R)$, aussi appelées \textbf{matrices de rotations}.
\end{notation}

\begin{prop}
	Les ensembles $\mathcal{O}_n(\R)$ et $\mathcal{SO}_n(\R)$ sont des sous-groupes de $\mathcal{GL}_n(\R)$.
\end{prop}

\begin{terminologie}
	Ces groupes sont respectivement appelés \textbf{groupe orthogonal d'ordre n} et \textbf{groupe spécial orthogonal d'ordre n}.
\end{terminologie}

\subsection{Espaces euclidiens orientés}

Dans cette sous-section, on suppose que $E$ est un espace euclidien non nul.

\begin{rappel}
	On définit sur l'ensemble des bases de $E$ une relation d'équivalence en posant :
$$\mathcal{B}\mathcal{R}\mathcal{B}'\iff \det_\mathcal{B}(\mathcal{B}')>0$$
\begin{itemize}
	\item Pour cette relation, il y a deux classes d'équivalence.
	\item Orienter l'espace $E$ revient à choisir l'une de ces deux classes et décréter "directes" ses éléments.
	\item On oriente $\R^n$, pour $n>0$, en décrétant directe sa base canonique.
	\item Lorsque $\mathcal{B}\mathcal{R}\mathcal{B}'$, on dit que les bases $\mathcal{B}$ et $\mathcal{B}'$ définissent la même orientation.
\end{itemize}
\end{rappel}

\begin{prop}
	Soit $\mathcal{B}$ une base orthonormée de $E$ et $\mathcal{B}'$ une famille de $n$ vecteurs de $E$. Alors $\mathcal{B}'$ est une base orthonormée de même orientation que $\mathcal{B}$ si, et seulement si, $M_\mathcal{B}(\mathcal{B}')\in\mathcal{SO}_n(\R)$.
\end{prop}

\begin{rmq}
	Les matrices orthogonales positives sont donc les matrices de changement de bases orthonormées directes. En particulier, $A\in\mathcal{SO}_n(\R)$ si, et seulement si, la famille des colonnes de $A$ est une base orthonormée directe de $\R^n$.
\end{rmq}

\begin{prop}
	Soit $\mathcal{B}$ et $\mathcal{B}'$ deux bases orthonormées directes de $E$. Alors on a : $$\det_\mathcal{B}=\det_{\mathcal{B}'}.$$
\end{prop}

\begin{terminologie}
	Pour $(u_1,\ldots,u_n)\in E^n$, le réel $\det_\mathcal{B}(u_1,\ldots,u_n)$, indépendant de la base orthonormée directe $\mathcal{B}$ choisie, est appelé \textbf{produit mixte} de $(u_1,\ldots,u_n)$ et parfois noté $[u_1,\ldots,u_n]$.
\end{terminologie}

\begin{dfn}
	On suppose $n\geqslant 2$. Pour $(u_1,\ldots,u_{n-1})\in E^{n-1}$, on considère $$\begin{array}{lrcl}
		\varphi:&E&\longrightarrow&\R \\
		&x&\longmapsto&[u_1,\ldots,u_{n-1},x]
	\end{array}$$ $\varphi$ est une forme linéaire sur $E$ donc il existe un unique $w\in E$ tel que $\forall x\in E,\ \varphi(x)=(w\mid x)$. $w$ est appelé le \textbf{produit vectoriel} de $u_1,\ldots,u_{n-1}$. On le note $u_1\wedge \cdots\wedge u_{n-1}$.
\end{dfn}

\begin{rmq}
	L'application : $$\begin{array}{rcl}
		E^{n-1}&\longrightarrow& E\\
		(u_1,\ldots,u_{n-1})&\longmapsto&u_1\wedge\cdots\wedge u_{n-1}
	\end{array}$$ est $(n-1)$-linéaire alternée.
\end{rmq}

\section{Endomorphismes autoadjoints}

\subsection{Généralités}

\begin{dfn}
	On dit qu'un endomorphisme $u\in\mathcal{L}(E)$ est \textbf{autoadjoint} lorsque $u\st=u$, c'est-à-dire lorsque $$\forall (x,y)\in E^2,\ (x\mid u(y))=(u(x)\mid y).$$
\end{dfn}

\begin{notation}
	L'ensemble des endomorphismes autoadjoints de $E$ est noté $\mathcal{S}(E)$.
\end{notation}

\begin{notation}
	Un endomorphisme autoadjoint est parfois appelé \textbf{endomorphisme symétrique}, ce qui explique la notation $\mathcal{S}(E)$.
\end{notation}

\begin{prop}
	L'ensemble $\mathcal{S}(E)$ est un sous-espace vectoriel de $\mathcal{L}(E)$.
\end{prop}

\begin{prop}
	Soit $\mathcal{B}$ une base orthonormée de $E$ et $u\in\mathcal{L}(E)$. Alors : $$u\in \mathcal{S}(E)\iff M_\mathcal{B}(u)\in\mathcal{S}_n(\R).$$
\end{prop}

\begin{rmq}
	On munit $\R^n$ de sa structure euclidienne canonique, si bien que la base canonique est une bas orthonormée. Soit $A\in\mathcal{M}_n(\R)$. L'endomorphisme canoniquement associé à la matrice $A$ est autoadjoint si, et seulement si, $A\in\mathcal{S}_n(\R)$.
\end{rmq}

\begin{attention}
	L'ensemble $\mathcal{S}(E)$ n'est en général pas stable par composition.
\end{attention}

\begin{cor}
	L'espace vectoriel $\mathcal{S}(E)$ est de dimension $\displaystyle\frac{n(n+1)}{2}$.
\end{cor}

\begin{lem}
	Soit $u\in\mathcal{S}(E)$ et $F$ un sous-espace vectoriel stable par $u$. Alors l'endomorphisme induit $u_F$ est autoadjoint.
\end{lem}

\subsubsection*{Projecteurs orthogonaux}

\begin{dfn}
	Soit $p$ un projecteur de $E$. On dit que $p$ est un \textbf{projecteur orthogonal} si $\ker(p)$ et $\text{Im}(p)$ sont orthogonaux.
\end{dfn}

\begin{rmq}
	Si $p$ est un projecteur orthogonal, on a la décomposition : $$E=\text{Im}(p)\overset{\bot}{\oplus}\ker(p)\ \ \text{donc} \ \ \text{Im}(p)=\ker(p)^\bot\ \ \text{et}\ \ \ker(p)=\text{Im}(p)^\bot$$
\end{rmq}

\begin{prop}
	Un projecteur de $E$ est orthogonal si, et seulement s'il est autoadjoint.
\end{prop}

\begin{rmq}
	On munit $\R^n$ de sa structure euclidienne canonique. Soit $A\in\mathcal{M}_n(\R)$. L'endomorphisme canoniquement associé à $A$ est un projecteur orthogonal si, et seulement si, $A^2=A$ et $A\in\mathcal{S}_n(\R)$.
\end{rmq}

\begin{exo}
	Soit $p$ un projecteur de $E$. Montrer que $p$ est un projecteur orthogonal si, et seulement si, $\forall x\in E,\ \lVert p(x)\rVert\leqslant\lVert x\rVert$.
\end{exo}

\subsection{Réduction des endomorphismes}

\begin{lem}
	Soit $u$ un endomorphisme d'un $\R$-espace vectoriel $E$ de dimension finie non nulle. Alors il existe une droite ou un plan stable par $u$.
\end{lem}

\begin{prop}
	Soit $u\in\mathcal{S}(E)$ et $F$ sous-espace vectoriel stable par $u$. Alors $F^\bot$ est stable par $u\st$.
\end{prop}

\begin{lem}
	Les sous-espaces propres d'un endomorphisme autoadjoint sont deux à deux orthogonaux.
\end{lem}

Dans le théorème suivant et dans le reste du chapitre, $E_\lambda(u)$ désigne le sous-espace propre de l'endomorphisme $u$ associé à la valeur propre $\lambda$.

\begin{thm}[Théorème spectral]
	Soit $u\in\mathcal{L}(E)$. les propositions suivantes sont équivalentes :
	\begin{enumerate}
		\item $u$ est autoadjoint ;
		\item $\displaystyle E=\bigoplus_{\lambda\in\text{Sp}(u)}^{\bot}E_\lambda(u)$ ;
		\item il existe une base orthonormée de $E$ diagonalisant $u$.
	\end{enumerate}
\end{thm}

\begin{ex}
	Si $u$ est un endomorphisme autoadjoint d'un espace vectoriel euclidien non nul, on a : $$\underset{\lVert x\rVert=1}{\min}(x\mid u(x))=\min\text{Sp}(u)\ \ \text{et}\ \ \underset{\lVert x\rVert=1}{\max}(x\mid u(x))=\max\text{Sp}(u)$$
\end{ex}

\begin{exo}
	\begin{enumerate}
		\item Montrer que pour tout $u\in\mathcal{S}(E)$, on a $\lVert u\rVert_{op}=\displaystyle\underset{\lambda\in\text{Sp}(u)}{\max}\lvert\lambda\rvert$.
		\item Montrer que pour tout $u\in\mathcal{L}(E)$, on a $\lVert u\rVert_{op}=\sqrt{\lVert u\st\circ u\rVert_{op}}=\displaystyle\underset{\lambda\in\text{Sp}(u\st\circ u)}{\max}\sqrt{\lambda}$.
	\end{enumerate}
\end{exo}

\begin{cor}
	Une matrice $A\in\mathcal{M}_n(\R)$ est symétrique si, et seulement si, elle est orthogonalement semblable à une matrice diagonale.
\end{cor}

\begin{rmq}
	Le corollaire précédent fournit une condition suffisant extrêmement simple de diagonalisabilité d'une matrice réelle : il suffit qu'elle soit symétrique.
\end{rmq}

\begin{attention}
	Ce résultat n'est plus vrai pour les matrices complexes. La matrice complexe symétrique $\displaystyle\begin{pmatrix}
	i&1\\
	1&-i
\end{pmatrix}$ est nilpotente, non nulle, donc n'est pas diagonalisable.
\end{attention}

\begin{ex}
	On considère $A=\displaystyle\begin{pmatrix}
		1&1&1\\
		1&1&1\\
		1&1&1
	\end{pmatrix}\in\mathcal{S}_3(\R)$.
	\begin{itemize}
		\item On a $\text{rg}(A)=1$ donc $0$ est valeur propre de multiplicité $2$ puisque $A$ est symétrique réelle donc diagonalisable. Le noyau de $A$ est le plan d'équation $x+y+z=0$ dont $(e_1,e_2)$ est une base orthonormée avec $$e_1=\frac{1}{\sqrt{2}}\begin{pmatrix}
			1\\
			-1\\
			0
		\end{pmatrix}\ \ \text{et}\ \ e_2=\frac{1}{\sqrt{6}}\begin{pmatrix}
		1\\
		1\\
		-2
		\end{pmatrix}$$
		\item Par ailleurs, $\Tr(A)=3$ donc $3$ est la dernière valeur propre et $e_3=\displaystyle\frac{1}{\sqrt{3}}\begin{pmatrix}
			1\\
			1\\
			1
		\end{pmatrix}$ est un vecteur propre unitaire associé.
		\item On pose $P\in\mathcal{O}_3(\R)$ la matrice dont les colonnes sont $(e_1,e_2,e_3)$ et, par les formules de changement de bases, on a : $$A=PDP\inv\ \ \text{avec}\ \ D=\text{Diag}(0,0,3)\ \ \text{et}\ \ P\inv=P^T$$
	\end{itemize}
\end{ex}

\begin{exo}
	Soit $A\in\mathcal{M}_n(\R)$. On suppose que pour tout matrice $P\in\mathcal{O}_n(\R)$, la matrice $P\inv AP$ est à diagonale nulle. Montrer que $A$ est antisymétrique.
\end{exo}

\begin{thm}[Représentation des formes bilinéaires (HP)]
	Pour tout endomorphisme $u$ de $E$, on note : $$\begin{array}{lrcl}
		\varphi_u:&E^2&\longrightarrow&\R \\
		&(x,y)&\longmapsto&(x\mid u(y))
	\end{array}$$
	\begin{enumerate}
		\item L'application $u\mapsto\varphi_u$ est un isomorphisme de $\mathcal{L}(E)$ sur $\mathcal{BL}(E)$.
		\item Elle induit un isomorphisme de $\mathcal{S}(E)$ sur $\mathcal{BLS}(E)$.
	\end{enumerate}
\end{thm}

\begin{exo}
	Soit $q$ une forme quadratique sur $E$. Montrer qu'il existe une base orthonormée $(e_1,\ldots,e_n)$ de $E$ et des scalaires ($\lambda_1,\ldots,\lambda_n)\in\R^n$ tels que $$\forall (x_1,\ldots,x_n)\in\R^n,\ q\left(\sum_{i=1}^{n}x_ie_i\right)=\sum_{i=1}^{n}\lambda_ix_i^2$$
\end{exo}

\begin{rmq}
	Soit $A\in\mathcal{S}_n(\R)$. $A$ représente une forme quadratique sur $\R^n$. Si on note $B$ la sous-matrice de $A$ obtenue en conservant les coefficients $(i,j)\in\llbracket1,n\rrbracket^2$, alors $B$ représente la restriction de la forme quadratique à $\R^{n-1}$.
\end{rmq}

\subsection{Endomorphismes autoadjoints (définis) positifs}

\begin{dfn}
	Soit $u\in\mathcal{S}(E)$.
	\begin{itemize}
		\item On dit que $u$ est autoadjoint \textbf{positif} lorsque $$\forall x\in E,\ (x\mid u(x))\geqslant 0$$
		\item On dit que $u$ est autoadjoint \textbf{défini positif} lorsque $$\forall x\in E\setminus\{0\},\ (x\mid u(x))> 0$$
	\end{itemize}
\end{dfn}

\begin{notation}
	L'ensemble des endomorphismes autoadjoints positifs (respectivement définis positifs) de $E$ est noté $\mathcal{S}^+(E)$ (respectivement $\mathcal{S}^{++}(E)$).
\end{notation}

\begin{rmq}
	Comme $u$ est autoadjoint, la forme bilinéaire $(x,y)\mapsto (x\mid u(y))$ est symétrique.
	\begin{itemize}
		\item Elle est positive si, et seulement si, $u$ est positif.
		\item C'est un produit scalaire si, et seulement si, $u$ est défini positif.
	\end{itemize}
\end{rmq}

\begin{ex}
	\begin{enumerate}
		\item Soit $u\in\mathcal{L}(E)$. Alors $u\st\circ u\in\mathcal{S}^+(E)$.
		\item L'exemple précédent montrer que $u$ est injectif (et donc bijectif en dimension finie) si, et seulement si, $u\st\circ u\in\mathcal{S}^{++}(E)$.
	\end{enumerate}
\end{ex}

La propriété suivante fournit une caractérisation des endomorphismes autoadjoints (définis) positifs à l'aide de leur spectre.

\begin{prop}
	Soit $u\in\mathcal{S}(E)$. On a les équivalences suivantes : 
	\begin{itemize}
		\item $u\in\mathcal{S}^+(E)\iff \text{Sp}(u)\subset\R_+$ ;
		\item $u\in\mathcal{S}^{++}(E)\iff \text{Sp}(u)\subset\R_+\st$.
	\end{itemize}
\end{prop}

\begin{ex}
	Si $u$ est un endomorphisme autoadjoint d'un espace vectoriel euclidien non nul, on a : $$\underset{\lVert x\rVert=1}{\min}(x\mid u(x))=\min\text{Sp}(u)\ \ \text{et}\ \ \underset{\lVert x\rVert=1}{\max}(x\mid u(x))=\max\text{Sp}(u)$$
\end{ex}

Munissons $\R^n$ de sa structure euclidienne canonique. En appliquant les définitions et caractérisations qui précèdent à l'endomorphisme canoniquement associé à une matrice, on obtient les définitions suivantes.

\begin{dfn}
	Soit $A\in\mathcal{S}_n(\R)$. On dit que $A$ est une matrice :
	\begin{itemize}
		\item \textbf{symétrique positive} lorsque : $$\forall x\in\R^n,\ (x\mid Ax)\geqslant 0\ \ \text{ou encore}\ \ \text{Sp}(A)\subset\R_+$$
		\item \textbf{symétrique définie positive} lorsque : $$\forall x\in\R^n\setminus\{0\},\ (x\mid Ax)\geqslant 0\ \ \text{ou encore}\ \ \text{Sp}(A)\subset\R_+\st$$
	\end{itemize}
\end{dfn}

\begin{notation}
	On note respectivement $\mathcal{S_n}^+(\R)$ et $\mathcal{S}_n^{++}(\R)$ l'ensemble des matrices symétriques positives et l'ensemble des matrices symétriques définies positives.
\end{notation}

\begin{ex}
	\begin{enumerate}
		\item La matrice nulle est une matrice symétrique positive.
		\item Par caractérisation par le spectre, si une matrice symétrique est définie positive, elle a une trace et un déterminant strictement positifs. La réciproque est fausse dès que $n\geqslant 3$, comme le montre l'exemple de la matrice $\text{Diag}(-1,-1,3,\ldots,3)$.
		\item Soit $A\in\mathcal{M}_n(\R)$. En notant $(e_1,\ldots,e_n)$ la base canonique de $\R^n$, on a : $$\forall i\in\llbracket1,n\rrbracket,\ a_{i,i}=(e_i\mid Ae_i)$$ Une matrice symétrique (définie) positive a donc tous ses coefficients diagonaux (strictement) positifs.
		\item Soit $A\in\mathcal{S}_2(\R)$. Alors : $$A\in\mathcal{S}_2^{++}(\R)\iff \Tr(A)>0\ \text{et}\ \det(A)>0$$
	\end{enumerate}
\end{ex}

\subsection{Applications (HP)}

\begin{thm}[Pseudo-réduction simultanée]
	Soit $A\in\mathcal{S}_n^{++}(\R)$ et $B\in\mathcal{S}_n(\R)$. Alors il existe une matrice $P\in\mathcal{GL}_n(\R)$ et une matrice diagonale $D\in\mathcal{M}_n(\R)$ telles que $$A=P^TP\ \ \text{et}\ \ B=P^TDP$$
\end{thm}

\begin{proof}
	On considère $E$ l'espace vectoriel euclidien $\R^n$ muni du produit scalaire $(X,Y)\mapsto X^TAY$. Ensuite, $q:X\mapsto X^TBX$ est une forme quadratique sur $E$. On prend alors $u\in\mathcal{S}(E)$ telle que $\forall x\in E,\ q(x)=(x\mid u(x))$ d'après le théorème de représentation des formes bilinéaires symétriques. On fixe $\mathcal{B}$ une base orthonormée de diagonalisation de $u$. Alors : $$M_\mathcal{B}(u)=D\ \ \text{donc} \ \ M_\mathcal{B}(q)=D$$ Or, on a aussi $M_{\mathcal{B}_c}(q)=B$ et on a : $$M_\mathcal{B}((\cdot\mid\cdot))=I_n \ \ \text{et} \ \ M_{\mathcal{B}_c}((\cdot\mid\cdot))=A$$ Par conséquent, il existe $P\in\mathcal{GL}_n(\R)$ telle que $$A=P^TI_nP\ \ \text{et}\ \ B=P^TDP$$
\end{proof}

\begin{thm}[Racine carrée]
	Soit $u$ un endomorphisme autoadjoint positif. Alors, il existe un unique endomorphisme autoadjoint positif $v$ tel que $u=v^2$.
\end{thm}

\begin{cor}[Décomposition polaire]
	Soit $M\in\mathcal{GL}_n(\R)$. Alors, il existe un unique couple $(O,S)\in\mathcal{O}_n(\R)\times\mathcal{S}_n^{+}(\R)$ tel que $M=OS$.
\end{cor}

\begin{exo}
	Montrer que le résultat d'existence subsiste pour tout matrice de $\mathcal{M}_n(\R)$. Y a-t-il encore unicité ?
\end{exo}

\section{Isométries vectorielles}

\subsection{Généralités}

\begin{dfn}
	Soit $u\in\mathcal{L}(E)$. On dit que $u$ est une \textbf{isométrie vectorielle} lorsque $$\forall x\in E,\ \lVert u(x)\rVert=\lVert x\rVert$$
\end{dfn}

\begin{notation}
	On note $\mathcal{O}(E)$ l'ensemble des isométries vectorielles de $E$.
\end{notation}

\begin{rmq}
	Attention ! La notion d'isométrie vectorielle dépend de la norme euclidienne et donc du produit scalaire. On pourra donc, dans certains cas où ce n'est pas clair, indexer $\mathcal{O}(E)$ par le produit scalaire correspondant.
\end{rmq}

\begin{prop}
	Soit $u\in\mathcal{L}(E)$. Soit $\mathcal{B}$ une base orthonormée de $E$. Les propriétés suivantes sont équivalentes : \begin{enumerate}[label=(\roman*)]
		\item $u\in\mathcal{O}(E)$ ;
		\item $\forall (x,y)\in E^2,\ \left(u(x)\mid u(y)\right)=(x\mid y)$ ;
		\item $u\in\mathcal{GL}(E)$ et $u\inv=u\st$ ;
		\item $M_\mathcal{B}(u)\in\mathcal{O}_n(\R)$ ;
		\item la famille $u(\mathcal{B})$ est une base orthonormée de $E$.
	\end{enumerate}
\end{prop}

\begin{terminologie}
	Une isométrie étant un élément de $\mathcal{GL}(E)$, on parle parfois d'\textbf{automorphisme orthogonal} pour désigner un tel élément, ce qui explique la notation $\mathcal{O}(E)$.
\end{terminologie}

\begin{prop}
	Soit $u\in\mathcal{O}(E)$. Alors, $\det(u)\in\{-1,1\}$.
\end{prop}

\begin{dfn}
	Une isométrie vectorielle $u\in\mathcal{O}(E)$ est dite :
	\begin{itemize}
		\item \textbf{positive} ou \textbf{directe} lorsque $u$ est de déterminant $1$ ;
		\item \textbf{négative} ou \textbf{indirecte} lorsque $u$ est de déterminant $-1$.
	\end{itemize}
\end{dfn}

\begin{notation}
	On note $\mathcal{SO}(E)$ l'ensemble des isométries positives de $E$, aussi appelées \textbf{rotations} de $E$.
\end{notation}

\begin{prop}
	Soit $u\in\mathcal{L}(E)$ et $\mathcal{B}$ une base orthonormée de $E$. Alors, on a : $$u\in\mathcal{SO}(E)\iff M_\mathcal{B}(u)\in\mathcal{SO}_n(\R).$$
\end{prop}

\begin{conséquence}
	On munit $\R^n$ de sa structure euclidienne orientée canonique. Soit $A\in\mathcal{M}_n(\R)$. La matrice $A$ est orthogonale (respectivement orthogonale positive) si, et seulement si, l'endomorphisme canoniquement associé à $A$ est une isométrie (respectivement une isométrie directe).
\end{conséquence}

\begin{cor}{}
	Les ensembles $\mathcal{O}(E)$ et $\mathcal{SO}(E)$ sont des sous-groupes de $\mathcal{GL}(E)$.
\end{cor}

\begin{terminologie}
	Les groupes $\mathcal{O}(E)$ et $\mathcal{SO}(E)$ sont appelés respectivement \textbf{groupe orthogonal de E} et \textbf{groupe spécial orthogonal de E}.
\end{terminologie}

\begin{prop}
	On suppose $E$ orienté et on se donne $u\in\mathcal{L}(E)$ et $\mathcal{B}$ une base orthonormée directe de $E$. Alors $u\in\mathcal{SO}(E)$ si, et seulement si, l'image de $\mathcal{B}$ par $u$ est une base orthonormée directe de $E$.
\end{prop}

\subsubsection*{Symétries orthogonales}

\begin{dfn}
	Soit $F$ un sous-espace vectoriel de $E$. On appelle \textbf{symétrie orthogonale} par rapport à $F$ la symétrie par rapport à $F$ parallèlement à $F^\bot$.
\end{dfn}

\begin{rmq}
	Pour cette définition, il suffit que $F$ soit de dimension finie, même si $E$ est de dimension infinie.
\end{rmq}

\begin{terminologie}
	Une symétrie orthogonale par rapport à un hyperplan est appelée \textbf{réflexion}.
\end{terminologie}

\begin{prop}
	Soit $s\in\mathcal{L}(E)$ une symétrie orthogonale. Alors $s\in\mathcal{O}(E)$.
\end{prop}

\begin{exo}
	\begin{enumerate}
		\item Soit $s\in\mathcal{L}(E)$ une réflexion par rapport à un hyperplan $H$. En considérant une base orthonormée $\mathcal{B}$ adaptée à la décomposition $E=H\oplus H^\bot$, on a : $$M_\mathcal{B}(s)=\begin{pmatrix}
			I_{n-1}&(0)\\
			(0)&-1
		\end{pmatrix}$$ donc $\det(s)=-1$. Par ailleurs, $s$ est une isométrie vectorielle puisque c'est une symétrie orthogonale. Une réflexion est donc une isométrie vectorielle indirecte.
		\item Soit $s$ une symétrie de $E$ vérifiant $s\in\mathcal{O}(E)$. Alors on a : $$s^2=\text{Id}_E\ \ \text{et}\ \ \ s\st\circ s=\text{Id}_E$$ donc $s=s\st$. Comme $s$ est un endomorphisme autoadjoint, les sous-espaces propres $\ker(s-\text{Id}_E)$ et $\ker(s+\text{Id}_E)$ sont orthogonaux, donc $s$ est une symétrie orthogonale.
	\end{enumerate}
\end{exo}

\subsection{Isométries en dimension $2$}

\begin{prop}
	Soit $A\in\mathcal{M}_2(\R)$. Alors, on dispose des caractérisations suivantes :
	\begin{itemize}
		\item $A\in\mathcal{SO}_2(\R)\iff \exists\theta\in\R\ : \ A=\displaystyle\begin{pmatrix}
			\cos(\theta)&-\sin(\theta)\\
			\sin(\theta)&\cos(\theta)
		\end{pmatrix}$
		\item $A\in\mathcal{O}_2(\R)\setminus\mathcal{SO}_2(\R)\iff \exists\theta\in\R\ : \ A=\displaystyle\begin{pmatrix}
			\cos(\theta)&\sin(\theta)\\
			\sin(\theta)&-\cos(\theta)
		\end{pmatrix}$
	\end{itemize}
\end{prop}

\begin{prop}
	L'application : $$\begin{array}{lrcl}
		R:&(\R,+)&\longrightarrow&(\mathcal{SO}_2(\R),\times)\\
		&\theta&\longmapsto&\displaystyle\begin{pmatrix}
			\cos(\theta)&-\sin(\theta)\\
			\sin(\theta)&\cos(\theta)
		\end{pmatrix} 
	\end{array}$$ est un morphisme de groupes surjectif de noyau $2\pi\Z$.
\end{prop}

\begin{cor}
	Le groupe $\mathcal{SO}_2(\R)$ est commutatif et isomorphe à $\U$.
\end{cor}

Dans le reste de cette section, $E$ désigne un plan euclidien orienté, c'est-à-dire un espace euclidien de dimension $2$ que l'on oriente par le choix d'une base orthonormée directe de référence.

\begin{prop}
	Soit $u\in\mathcal{SO}(E)$. Il existe $\theta\in\R$, unique modulo $2\pi$, tel que pour toute base orthonormée directe $\mathcal{B}$, on ait : $$M_\mathcal{B}(u)=\begin{pmatrix}
		\cos(\theta)&-\sin(\theta)\\
		\sin(\theta)&\cos(\theta)
	\end{pmatrix}$$
\end{prop}

\begin{terminologie}
	Un tel élément de $\mathcal{SO}(E)$ est appelé \textbf{rotation d'angle $\theta$}. Le réel $\theta$ est appelé la \textbf{mesure de l'angle de la rotation}.
\end{terminologie}

\begin{rmq}
	Si $\mathcal{B}=(e_1,e_2)$ est une base orthonormée indirecte de $E$, et $u$ une rotation d'angle $\theta$, alors $\mathcal{B}'=(-e_1,e_2)$ est une base orthonormée directe de $E$, donc on a $M_{\mathcal{B}'}(u)=R(\theta)$ $M_{\mathcal{B}}(u)=R(-\theta)$ par une vérification immédiate.
\end{rmq}

\begin{cor}
	Le groupe $\mathcal{SO}(E)$ est commutatif et isomorphe à $\U$.
\end{cor}

\begin{rmq}
	Soit $(u,v)\in\mathcal{SO}(E)^2$ deux rotations d'angles respectifs $\theta$ et $\theta'$. Si $\mathcal{B}$ est une base orthonormée directe de $E$, on a : $$M_\mathcal{B}(u\circ v)=M_\mathcal{B}(u)M_\mathcal{B}(v)=R(\theta)R(\theta')=R(\theta+\theta')$$ donc $u\circ v$ est la rotation d'angle $\theta+\theta'$.
\end{rmq}

\begin{prop}
	Soit $u\in\mathcal{O}(E)\setminus\mathcal{SO}(E)$. Alors $u$ est une réflexion de $E$.
\end{prop}

\subsubsection*{Angles orientés de vecteurs dans $E$}

\begin{prop}
	Soit $x$ et $y$ deux vecteurs unitaires de $E$. Il existe une unique rotation $u\in\mathcal{SO}(E)$ telle que $u(x)=y$.
\end{prop}

\begin{dfn}
	Soit $x$ et $y$ deux vecteurs non nuls du plan $E$.
	\begin{itemize}
		\item Il existe une unique rotation $u$ qui envoie le renormalisé de $x$ sur le renormalisé de $y$.
		\item L'angle de cette rotation est appelé \textbf{angle orienté des vecteurs x et y}.
	\end{itemize}
\end{dfn}

\begin{notation}
	L'angle orienté des vecteurs non nuls $x$ et $y$ est noté $\widehat{(x,y)}$.
\end{notation}

\begin{prop}
	Soit $(x,y,z)\in E^3$ trois vecteurs non nuls. Alors, on a : $$(x,z)=(x,y)+(y,z)\ [2\pi].$$
\end{prop}

\subsection{Réduction des isométries vectorielles}

\begin{lem}
	L'endomorphisme induit sur un sous-espace vectoriel stable par une isométrie vectorielle est aussi une isométrie vectorielle.
\end{lem}

\begin{prop}
	Soit $u$ une isométrie vectorielle et $F$ un sous-espace vectoriel de $E$. Si $F$ est stable par $u$, alors $F^\bot$ est aussi stable par $u$.
\end{prop}

\begin{thm}[Réduction des isométries vectorielles]
	Soit $u$ une isométrie vectorielle. Il existe une base orthonormée $\mathcal{B}$ de $E$ dans laquelle la matrice de $u$ est égale à une matrice diagonale par blocs avec des blocs diagonaux :
	\begin{itemize}
		\item de taille $1$ et de la forme $(\gamma)$ avec $\gamma\in\{-1,1\}$ ;
		\item de taille $2$ et de la forme : $$R(\theta)=\begin{pmatrix}
			\cos(\theta)&-\sin(\theta)\\
			\sin(\theta)&\cos(\theta)
		\end{pmatrix}\ \ \text{avec}\ \ \theta\in\R\setminus\pi\Z$$
	\end{itemize}
\end{thm}

Autrement dit, la matrice d'une isométrie dans une telle base orthonormée peut se mettre sous la forme : $$\begin{pmatrix}
	I_p&&&&\\
	&-I_q&&(0)&\\
	&&R(\theta_1)&&\\
	&(0)&&\ddots&\\
	&&&&R(\theta_r)
\end{pmatrix}\ \ \text{avec}\ \ (\theta_1,\ldots,\theta_r)\in\left(\R\setminus\pi\Z\right)^r$$

\begin{exo}
	Soit $A\in\mathcal{S}_n(\R)$. Montrer que les deux propriétés suivantes sont équivalentes : 
	\begin{enumerate}[label=(\roman*)]
		\item $\exists B\in\mathcal{O}_n(\R)\ : \ A=B+B^T$ ;
		\item $\text{Sp}(A)\subset[-2,2]$ et les valeurs propres de $A$ dans $]-2,2[$ sont de multiplicité paire.
	\end{enumerate}
\end{exo}

\begin{conséquence}
	Toute matrice orthogonale est orthogonalement semblable à une matrice de la forme précédente.
\end{conséquence}

\begin{exo}
	Soit $u\in\mathcal{SO}(E)$. Montrer qu'il existe une base orthonormée de $E$ dans laquelle la matrice de $u$ est de la forme : $$\begin{pmatrix}
		I_s&&&(0)\\
		&R(\varphi_1)&&\\
		&&\ddots&\\
		(0)&&&R(\varphi_t)
	\end{pmatrix}\ \ \text{avec}\ \ s\in\{0,1\}\ \ \text{et}\ \ (\varphi_1,\ldots,\varphi_t)\in\R^t$$
\end{exo}

\begin{exo}
	\begin{enumerate}
		\item Montrer que $\mathcal{SO}_n(\R)$ est connexe par arcs. 
		\item En utilisant la décomposition polaire, en déduire que $\mathcal{GL}_n^+(\R)$ est connexe par arcs.
	\end{enumerate}
\end{exo}

\subsubsection*{Rotations en dimension 3}
Dans le reste de cette section, $E$ désigne un espace euclidien de dimension $3$ que l'on oriente par le choix d'une base orthonormée directe de référence.

\begin{prop}
	Si $u\in\mathcal{SO}(E)$, alors il existe une base orthonormée $\mathcal{B}$ de $E$ et un réel $\theta\in\R$ tels que : $$M_\mathcal{B}(u)=\begin{pmatrix}
		1&0&0\\
		0&\cos(\theta)&-\sin(\theta)\\
		0&\sin(\theta)&\cos(\theta)
	\end{pmatrix}$$
\end{prop}

\begin{conséquence}
	On munit $\R^3$ de sa structure euclidienne orientée canonique et on se donne $A\in\mathcal{SO}_3(\R)$. Alors il existe $\theta\in\R$ tel que $A$ soit orthogonalement semblable à une matrice de la forme $$\begin{pmatrix}
	1&0&0\\
	0&\cos(\theta)&-\sin(\theta)\\
	0&\sin(\theta)&\cos(\theta)
\end{pmatrix}$$
\end{conséquence}

\begin{cor}
	Si $u$ est une isométrie vectorielle directe de $E$, il existe une droite $D$ et un plan $P$ orthogonal à $D$ tel que $u_{\mid D}=\text{Id}_{D}$ et tel que $u_{\mid P}$ soit une rotation.
\end{cor}

\begin{rmq}
	\begin{enumerate}
		\item On dit qu'une telle isométrie vectorielle directe est une rotation d'angle $\theta$. L'angle $\theta$ de la rotation peut être déterminé, au signe près, en utilisant le fait que $\Tr(u)=1+2\cos(\theta)$.
		\item Si $u\in\mathcal{O}(E)\setminus\mathcal{SO}(E)$ est une isométrie indirecte., alors $-u\in\mathcal{SO}(E)$. En effet, $u$ conserve la norme donc $-u$ également, et par ailleurs, puisque l'on travaille en dimension $3$, on a $\det(-u)=(-1)^3\det(u)=1$. Il existe donc une base orthonormée $\mathcal{B}$ et un réel $\theta$ tels que $$M_\mathcal{B}(-u)=\begin{pmatrix}
			1&0\\
			0&R(\theta)
		\end{pmatrix} \ \ \text{donc}\ \ M_\mathcal{B}(u)=\begin{pmatrix}
		-1&0\\
		0&R(\theta+\pi)
		\end{pmatrix}$$
	\end{enumerate}
\end{rmq}

\end{document}